% ৮.২ Interpreter
\section{Interpreter}

\begin{learningbox}
এই টপিক শেষে শিক্ষার্থী interpreter-এর কাজ এবং compiler থেকে পার্থক্য ব্যাখ্যা করতে পারবে।
\end{learningbox}

\begin{definitionbox}{Interpreter}
যে translator source program line by line translate ও execute করে এবং error পেলে সাধারণত সেই line-এ execution থামায়, তাকে interpreter বলা হয়।
\end{definitionbox}

\begin{center}
\begin{tabular}{|p{0.24\textwidth}|p{0.30\textwidth}|p{0.30\textwidth}|}
\hline
\textbf{ভিত্তি} & \textbf{Compiler} & \textbf{Interpreter} \\
\hline
Translation & পুরো program একসাথে & line by line \\
\hline
Execution speed & দ্রুত & তুলনামূলক ধীর \\
\hline
Error handling & error list একসাথে & line ধরে error দেখায় \\
\hline
Object code & তৈরি হতে পারে & সাধারণত তৈরি করে না \\
\hline
\end{tabular}
\end{center}

\begin{examplebox}{শেখার সুবিধা}
Interpreter line ধরে error দেখায় বলে beginners দ্রুত ভুল খুঁজতে পারে।
\end{examplebox}
